\documentclass[12pt]{article}
\usepackage[T1]{fontenc}
\usepackage[utf8]{inputenc}
\usepackage{lmodern}
\usepackage{textcomp}
\usepackage{enumitem}
\usepackage{geometry}
\usepackage{graphicx}
\usepackage{amsmath, amssymb}
\geometry{margin=1in}
\begin{document}
\begin{center}
Constitutional Law - Graduate Level Assessment 24 \
Total Marks: 40 \
Answer all questions.
\end{center}

\section*{Multiple Choice (10 marks)}
\begin{enumerate}[label=\arabic*.]
\item Which statement best describes one of Dworkin's central arguments in Justice for Hedgehogs?
\begin{enumerate}[label=(\alph*)]
    \item The law is the only guide for ethical behaviour.
    \item The law dictates what moral values should affect our ethical behaviour.
    \item Morality plays no role in the concept of law.
    \item Moral arguments operate only in hard cases.
    \item Morality only applies to specific, complex cases.
\end{enumerate}
\item Which case was the first to define the meaning of the doctrine of 'margin of appreciation' as applied by the European Court of Human Rights?
\begin{enumerate}[label=(\alph*)]
    \item Dudgeon v UK (1981)
    \item Marckx v Belgium (1979)
    \item Osman v UK ( 1998)
    \item Tyrer v UK (1978)
    \item McCann v UK (1995)
\end{enumerate}
\item Do treaties bind third States, ie non-State parties?
\begin{enumerate}[label=(\alph*)]
    \item Treaties do no create obligations or rights for third States without their consent
    \item Treaties create both obligations and rights for third States
    \item Treaties do not create any rights for third States, even when the latter consent.
    \item Treaties may create only obligations for third States
    \item Treaties create obligations for third States without their consent
\end{enumerate}
\item What is it called when a remainder in the grantor's heirs is invalid and becomes a reversion in the grantor?
\begin{enumerate}[label=(\alph*)]
    \item Related Doctrine of Merger
    \item Contingent remainder
    \item Doctrine of Worthier Title
    \item Doctrine of Escheat
    \item Vested remainder
\end{enumerate}
\item Maine's famous aphorism that 'the movement of progressive societies has hitherto been a movement from Status to Contract' is often misunderstood. In what way?
\begin{enumerate}[label=(\alph*)]
    \item His concept of status is misrepresented.
    \item It is misinterpreted as a prediction.
    \item His idea is considered inapplicable to Western legal systems.
    \item It is incorrectly related to economic progression.
    \item It is taken literally.
\end{enumerate}
\end{enumerate}

\section*{True / False (10 marks)}
\textit{Answer True or False}
\begin{enumerate}[label=\arabic*.]
\setcounter{enumi}{5}
\item True or False: A warranty of merchantability guarantees that a product will meet all the buyer's specific needs without requiring additional advice.
\item True or False: The original position (POP) would prioritize wealth over the establishment of a compassionate society.
\item True or False: Passive personality jurisdiction is based on the relationship between the victim and the offender.
\item True or False: Bentham argued that law is solely a product of moral principles dictated by society's needs.
\item True or False: In some jurisdictions, solicitation is defined as a conspiracy to commit a crime, requiring an agreement between multiple parties.
\end{enumerate}

\section*{Long Form Response (20 marks)}
\textit{Answer each question with a clear and complete explanation.}
\begin{enumerate}[label=\arabic*.]
\setcounter{enumi}{10}
\item Critically analyze the strongest argument against ethical relativism's stance on human rights, discussing the role of cognitivism in this context.
\item Discuss the legal implications of bid retraction in auction settings, particularly in the context of Cross \& Black's liquidation, and analyze how this principle affects bidders' rights.
\end{enumerate}

\vfill
\noindent\textit{End of Paper}
\end{document}